\documentclass[11pt,a4paper]{article}
\usepackage[T1]{fontenc}
\usepackage[utf8]{inputenc}
\usepackage{lmodern,amsmath,amssymb,amsthm,mathtools}
\usepackage[textwidth=145mm,textheight=225mm,centering]{geometry}
\usepackage{microtype,booktabs,enumitem}
\usepackage[authoryear,round]{natbib}
\usepackage[colorlinks=true,linkcolor=blue,citecolor=blue,urlcolor=blue,breaklinks=true]{hyperref}
\usepackage{xurl}
\newcommand{\workbenchurl}{https://github.com/KokunoYumeto/algebraic-cycles-and-motives-workbench}
\newcommand{\recordurl}{https://github.com/KokunoYumeto/mathematics-workbenches}
\newcommand{\OO}{\mathcal O}
\newcommand{\PP}{\mathbf P}
\newcommand{\cris}{\mathrm{cris}}
\newcommand{\dR}{\mathrm{dR}}
\newcommand{\et}{\mathrm{\acute et}}
\newcommand{\Sym}{\operatorname{Sym}}
\newcommand{\Spec}{\operatorname{Spec}}
\newcommand{\coker}{\operatorname{coker}}
\newcommand{\rank}{\operatorname{rank}}
\newcommand{\Tor}{\operatorname{Tor}}
\newcommand{\HH}{\mathbb H}
\newtheorem{theorem}{Theorem}[section]
\newtheorem{proposition}[theorem]{Proposition}
\newtheorem{lemma}[theorem]{Lemma}
\newtheorem{corollary}[theorem]{Corollary}
\theoremstyle{definition}\newtheorem{definition}[theorem]{Definition}
\theoremstyle{remark}\newtheorem{remark}[theorem]{Remark}
\numberwithin{equation}{section}
\setlist[enumerate]{itemsep=3pt,topsep=5pt}
\emergencystretch=1em
\title{Counterexamples to integral weak Lefschetz in Grothendieck's \emph{Cristaux}, after Raynaud and Mukai}
\author{}
\date{23 November 2026}
\begin{document}
\maketitle
\begin{abstract}
Grothendieck's \emph{Cristaux} admits an integral crystalline weak Lefschetz assertion while studying unit Frobenius eigenvalues in vanishing cohomology. Starting from the surfaces of Raynaud and the exact-differential construction of Tango and Mukai, we construct a smooth projective surface with a very ample line bundle whose negative first power has one-dimensional first cohomology and whose higher negative powers have none. In characteristic two the underlying curve has genus seven, and the construction uses six sections and 556 ninth roots. A product with the projective line gives smooth threefold hyperplane sections for which integral second-crystalline restriction is injective but its cokernel has torsion. For the same embedding, every smooth section of degree at least two has torsion-free cokernel and identifies the full torsion subgroups. A second, iterated Tango threefold has noninjective integral restriction precisely in degrees one and two, with injectivity in every degree at least three. The surface construction also gives a counterexample in every positive characteristic, using crystalline base change and rational weak Lefschetz as published inputs. The mathematics and text were produced and checked by GPT-6 Astra Pro, including separate statement-only reconstructions of the counterexample constructions by the same model; no human has verified them.
\end{abstract}

\section{The assertion and the counterexample}

The purpose of this paper is to distinguish the integral weak Lefschetz assertion admitted in Grothendieck's manuscript from the rational theorem that remains valid. Let $k$ be an algebraically closed field of characteristic $p>0$, let $W=W(k)$, and write
\[
 C^j(V)=H^j_{\cris}(V/W),\qquad T^j(V)=C^j(V)_{\mathrm{tors}}
\]
for a smooth projective $k$-variety $V$. All torsion in these finite $W$-modules is $p$-primary. For a smooth hyperplane section $Y$ of a smooth projective $n$-fold $V$, set
\[
 E(Y)=\coker\{C^{n-1}(V)\longrightarrow C^{n-1}(Y)\}.
\]

In \emph{Torsion en cohomologie cristalline}, cote 4 of the Montpellier archive, Grothendieck writes in \S8.3:
\begin{quote}
``Nous admettrons que ces relations restent valables dans le cas de la cohomologie cristalline.''
\end{quote}
Here the admitted relations are that $E(Y)$ is torsion-free and that restriction identifies $T^{n-1}(V)$ and $T^{n-1}(Y)$; equation (8.3.3) also asserts injectivity. The quotation and the torsion assertion occur on archive sheet 3, and (8.3.3) on sheet 4, respectively PDF images 4 and 5. On archive sheet 2, the typed qualification imposing sufficiently large degree is crossed out in the sentence introducing the hyperplane section \citep[\S\S8.1, 8.3, archive sheets 2--4]{Grothendieck}. The examples below disprove the unrestricted torsion and injectivity assertions. They concern an admitted assertion in a manuscript, rather than a theorem Grothendieck claimed to have proved.

Here is the complete construction used in the main result. The details that make it smooth, very ample, and numerically explicit are proved in Sections~\ref{sec:surface}--\ref{sec:differentials}.

\begin{definition}[The characteristic-two surface]\label{def:explicit}
Let $C$ be the smooth projective curve with affine equation $y^2+y=x^{15}$, let $P_\infty$ be its point at infinity, and put $N=\OO_C(2P_\infty)$ and $T=N^3$. The exact differential $dx$ defines a class
\[
 0\ne\alpha\in\ker\{H^1(C,T^{-1})\xrightarrow{F}H^1(C,T^{-2})\}.
\]
Let $0\to\OO_C\to V\to T\to0$ be its extension. The projective bundle $\pi:P=\PP_C(V)\to C$ parametrizes one-dimensional quotients. Let $E_0$ be the section from $V\to T$ and let $G$ be the disjoint inseparable double multisection obtained from a splitting of $F_C^*V$. Thus
\[
 \OO_P(E_0)=\OO_P(1),\quad
 \OO_P(G)=\OO_P(2)\otimes\pi^*T^{-2}.
\]
With $M=\OO_P(-1)\otimes\pi^*N^2$, form the cubic cover
\[
 \psi:S=\Spec_P(\OO_P\oplus M\oplus M^2)\longrightarrow P
\]
whose multiplication is defined by the section of $M^{-3}$ vanishing on $E_0+G$. Write $E$ for the reduced inverse image of $E_0$, $f=\pi\psi$, and
\[
 L=\OO_S(2E)\otimes f^*N.
\]
Choose a general basis $b_1,\ldots,b_{556}$ of $H^0(S,L^9)$ and form
\[
 \rho:X=\Spec_S\left(\bigotimes_{i=1}^{556}
                  \bigoplus_{a=0}^{8}L^{-a}\right)\longrightarrow S,
 \qquad z_i^9=b_i.
\]
Finally put $A=\rho^*L$. General means that the divisors of the $b_i$ are smooth, pairwise transverse, and have no triple intersection. These conditions specify a nonempty open set of choices.
\end{definition}

\begin{theorem}[The explicit surface and integral threefold]\label{thm:main}
In characteristic two the construction in Definition~\ref{def:explicit} has the following properties.
\begin{enumerate}
\item The curve $C$ has genus $7$, and
\[
 h^0(C,N)=2,\quad h^0(C,N^3)=4,\quad
 h^0(S,L)=6,\quad h^0(S,L^9)=556.
\]
The surface $X$ is smooth, projective, and connected. The line bundle $A$ is very ample; its six pullback sections and 556 root sections embed $X$ in $\PP^{561}$. The degree of $\rho$ is $9^{556}$.
\item For every smooth $Y_d\in|A^d|$,
\[
 \dim_k\ker\{H^1(X,\OO_X)\to H^1(Y_d,\OO_{Y_d})\}
 =\begin{cases}1,&d=1,\\0,&d\ge2.\end{cases}
\]
\item Put $m=591\cdot9^{556}+2$, $Z=X\times\PP^1$, and
$B=A\boxtimes\OO_{\PP^1}(m)$. This is a very ample line bundle. For every smooth $Y_d\in|B^d|$, restriction
\[
 C^2(Z)\longrightarrow C^2(Y_d)
\]
is injective. Its cokernel is torsion-free if and only if $d\ge2$, and its map $T^2(Z)\to T^2(Y_d)$ is an isomorphism if and only if $d\ge2$. More precisely, in degree one,
\[
 \dim_k C^2(Y_1)[2]\ge\dim_k C^2(Z)[2]+1.
\]
\end{enumerate}
\end{theorem}

In words, smooth hyperplane sections can acquire crystalline torsion that is absent from the ambient variety's image, even though the integral restriction map remains injective. Increasing the degree to two repairs both torsion assertions for this fixed embedding. Thus ampleness and smoothness alone do not give the integral saturation that Grothendieck's heuristic argument admitted. The example leaves rational crystalline weak Lefschetz intact.

\begin{theorem}[Every characteristic]\label{thm:all-primes}
For every prime $p$ and every algebraically closed field $k$ of characteristic $p$, there are a smooth projective connected surface $X$ and a very ample line bundle $A$ with $H^1(X,A^{-1})\ne0$. There are also a smooth projective connected threefold $Z$ with very ample $B$ such that every smooth $Y\in|B|$ has
\[
 \dim_k C^2(Y)[p]>\dim_k C^2(Z)[p].
\]
In particular restriction on degree-two torsion is not surjective and its full crystalline cokernel has torsion. An explicit construction and sufficient degree bound $d\ge p^2-p$ for saturation in these product examples are given below.
\end{theorem}

\begin{theorem}[Failure of integral injectivity]\label{thm:intro-injectivity}
In characteristic two, let $(U,H)$ be the smooth threefold and very ample line bundle constructed in Definition~\ref{def:threefold}: two Tango iterations of degrees 21 and 5 start from $y^2+y=x^{843}$, and 843 simultaneous 445th roots make the polarization very ample. For every smooth $Y_d\in|H^d|$,
\[
 C^2(U)\longrightarrow C^2(Y_d)
 \quad\hbox{is injective exactly when }d\ge3.
\]
In degrees one and two its kernel killed by two has dimension at least 361 and one, respectively. The coherent restriction kernel in $H^1(\OO)$ has dimension 1203 in degree one, one in degree two, and zero in every higher degree.
\end{theorem}

Thus integral weak Lefschetz can fail by a kernel as well as by a nonsaturated image. These failures require separate examples and separate arguments. Theorem~\ref{thm:regularity} gives an explicit sufficient degree bound for every smooth projective embedding. The exact cokernel and torsion-group ranges for the second threefold are not determined here.

\paragraph{Credit and relation to earlier results.}
Raynaud constructed smooth polarized surfaces violating Kodaira vanishing \citep{Raynaud}; Mukai extended the exact-differential and iteration methods \citep[Theorems 1--2, pp.~1--2; Theorem 1.2, p.~3 of the preprint]{Mukai}. Zheng gives the cyclic surface construction and the applicable nonvanishing range \citep[\S2.2, pp.~3--4; Theorem 3.7, p.~8]{Zheng}. Our polarization is the endpoint $a=p$, $b=1$ of Zheng's range. Its nonvanishing is therefore a known input, and we give its direct computation because global generation and exact negative powers are essential here.

Coherent and flat-cohomology Lefschetz failures are also established in earlier work. Cotner and Zavyalov's published Example 3.3.6 derives a noninjective $\alpha_p$-restriction from a reduced connected ample divisor with $H^1(\OO(-Y))\ne0$; their Theorem 3.3.5 requires strong ampleness. Their footnote 4 credits Burt Totaro with observing that higher-dimensional examples of Mukai satisfy their hypotheses \citep[p.~1862]{CotnerZavyalov}. The point of the present construction is a specified very ample polarization and the exact crystalline degree threshold. No claim of priority for the general failure of Kodaira or coherent Lefschetz vanishing is made. Rational crystalline weak Lefschetz is used in the precise form of Nakkajima's Theorem 9.1 and Corollary 9.3 \citep[pp.~63--64]{NakkajimaYobuko}. Integral theorems with sufficient ampleness or lifting hypotheses have different hypotheses from our examples.

A recent neighboring account writes out Alpöge's sketch of Deligne's Frobenius eigenvalue bound and reports production with OpenAI's GPT-6 Astra \citep{Hexagon}. Its stated inputs include rational étale weak Lefschetz, and its conclusion concerns weights. The distinction relevant here is that rational eigenvalue bounds leave integral torsion invisible. Our examples concern that additional integral assertion in Grothendieck's discussion; they do not challenge the rational weight theorem.

\paragraph{Why the proof works.}
An ample line bundle with nonzero negative cohomology need not give a smooth hyperplane section. The first step is therefore to choose a Raynaud polarization that is generated. Simultaneously adjoining roots of a basis of a very ample power makes its pullback very ample while preserving the negative cohomology as a direct summand. The curve's pole semigroup makes the characteristic-two numbers exact. On the product with $\PP^1$, the negative cohomology becomes additional $H^1(\OO)$ on every smooth hyperplane. The conjugate spectral sequence and crystalline universal coefficients force more second-crystalline torsion. A separate relative de Rham calculation, including its differential, proves the sharp positive result in higher degrees. For the second threefold, two Tango iterations produce a much larger coherent kernel, and the root cover preserves the first Betti number. This forces a nonzero integral kernel. The exact degree-two conclusion also needs an integral identification with the base curve's first crystalline lattice; negative one-forms then control injectivity from degree three onward.

\section{A generated Raynaud polarization}\label{sec:surface}

This section constructs the surface and computes the cohomology needed for the hyperplane example. We include the cyclic-cover formulas with their projective-bundle convention.

\begin{lemma}[The curve]\label{lem:curve}
Put $\ell=p+1$ and
\[
 h=\begin{cases}8,&p=2,\\p^2-2,&p\text{ odd},\end{cases}
 \quad m_0=hp-1,\quad d_0=(p-1)h-2.
\]
The smooth completion of $y^p-y=x^{m_0}$ has one point $P_\infty$ at infinity, and
\[
 (dx)=pd_0P_\infty,\qquad 2g-2=pd_0.
\]
The integer $n=d_0/\ell$ is $2$ for $p=2$ and $p(p-2)$ for odd $p$. The line bundle $N=\OO_C(nP_\infty)$ is generated, and $T=N^\ell$ satisfies $T^p\simeq\Omega_C^1$.
\end{lemma}
\begin{proof}
The Artin--Schreier polynomial is irreducible over $k(x)$. Indeed, a rational solution of $a^p-a=x^{m_0}$ could have no finite pole, while its pole order at infinity after applying $a\mapsto a^p-a$ would be divisible by $p$. Since $p\nmid m_0$, no solution exists. A degree-$p$ Artin--Schreier polynomial either has a root or is irreducible. The affine curve is smooth and étale over the $x$-line.

At a point above infinity, valuations in the equation give $p\,v(y)=m_0v(x)$. The ramification index is therefore divisible by $p$; it equals $p$, there is a unique point, and $v(x)=-p$, $v(y)=-m_0$. Put $w=y/x^h$, $v=1/x$, and $c=h(p-1)$. The equation becomes
\[
 v+wv^c=w^p.
\]
Its derivative in $v$ is a unit at the origin, so $w$ is a uniformizer and $v$ has order $p$. Differentiation gives
\[
 dx=\frac{v^{c-2}}{1+cwv^{c-1}}\,dw.
\]
Thus $dx$ has order $p(c-2)=pd_0$ at infinity. It has no zero or pole elsewhere. Taking the degree of its divisor gives the genus formula. Direct substitution gives the asserted values of $n$. The sections $1$ and $x^{n/p}$ generate $N$: the first is nonzero away from infinity and the second has exactly the allowed pole order there. Finally the divisor of $dx$ gives $\Omega_C^1\simeq T^p$.
\end{proof}

The exact-differential sequence identifies $dx$ with a nonzero section of $B_C^1\otimes T^{-1}$, where $B_C^1=\operatorname{im}(F_{C*}d)$. In
\[
 0\longrightarrow T^{-1}\longrightarrow F_{C*}T^{-p}
                 \longrightarrow B_C^1\otimes T^{-1}\longrightarrow0,
\]
the middle term has no global sections. The connecting class $\alpha$ is consequently nonzero and killed by Frobenius. This is the exact-differential description of Tango's theorem \citep[Theorem 1.2 and Lemma 1.4, pp.~3--4 of the preprint]{Mukai}.

Let $0\to\OO_C\to V\to T\to0$ represent $\alpha$. On $P=\PP_C(V)$ let $E_0$ be its quotient section. A splitting of $F_C^*V$ gives a section of $\PP_C(F_C^*V)$ disjoint from the Frobenius image of $E_0$; its inverse image under relative Frobenius is a multisection $G$. The projective-bundle divisor formulas give
\begin{equation}\label{eq:branch}
 \OO_P(E_0)=\OO_P(1),\qquad
 \OO_P(G)=\OO_P(p)\otimes\pi^*T^{-p}.
\end{equation}
Locally away from $E_0$, the equation is $z^p=f_i$. The differential $df_i$ is $dx$ divided by a local generator of $T^p$, and is nowhere zero by Lemma~\ref{lem:curve}. The Jacobian criterion shows that $G$ is smooth. This explains the smoothness in the Tango--Raynaud construction without identifying purely inseparable maps with smooth maps.

Set
\[
 M=\OO_P(-1)\otimes\pi^*N^p,
 \qquad M^{-\ell}=\OO_P(E_0+G).
\]
Let $\psi:S=\Spec_P(\bigoplus_{i=0}^{p}M^i)\to P$ be the cyclic cover with this branch divisor. Since $E_0$ and $G$ are disjoint and smooth and $\ell$ is invertible, its local equations give a smooth surface. The valuation at $E_0$ makes the generic equation Eisenstein, so $S$ is integral. Write $E$ for the reduced inverse image of $E_0$, and $f=\pi\psi$. Then $\psi^*E_0=\ell E$. The map $E\to C$ is an isomorphism and $E^2=n$.

\begin{lemma}[Cyclic eigensheaves]\label{lem:eigen}
For every integer $a$,
\begin{equation}\label{eq:eigen}
 \psi_*\OO_S(aE)=
 \bigoplus_{i=0}^{p}\OO_P\!\left(\left\lfloor\frac{a-pi}{\ell}\right\rfloor\right)
                         \otimes\pi^*N^{pi}.
\end{equation}
\end{lemma}
\begin{proof}
At $E_0$, choose a branch parameter $s$ and a cyclic coordinate $u$ with $u^\ell=s$ times a unit. The $i$th character has generator $u^i$. Allowing poles of order at most $a$ on $E$ allows the coefficient of $u^i$ a pole of order $\lfloor(a+i)/\ell\rfloor$ on $E_0$. Hence its eigensheaf is
$M^i(\lfloor(a+i)/\ell\rfloor E_0)$. There is no additional pole at $G$ or elsewhere. Since $M^i=\OO_P(-i)\otimes\pi^*N^{pi}$ and $\ell=p+1$, this is precisely the displayed formula. Regularity can be checked in codimension one because both the surface and the indicated locally free sheaves are normal.
\end{proof}

\begin{proposition}[The polarization]\label{prop:polarization}
The line bundle $L=\OO_S(pE)\otimes f^*N$ is ample and generated, $H^1(S,L^{-1})\ne0$, and $L^{\ell^2}$ is very ample. Moreover
\begin{equation}\label{eq:pushL}
 f_*L=N\oplus N^{p+1}.
\end{equation}
\end{proposition}
\begin{proof}
Each geometric fibre of $f$ is integral: its affine equation is $u^\ell=(z-a)^p$, irreducible because $\gcd(\ell,p)=1$, with one point at infinity. The section $E$ has degree one on it. The effective divisor $E$ is nef, since $E^2=n>0$ and its intersection with every different irreducible curve is nonnegative. On a vertical curve $pE$ has positive degree, and on a horizontal curve $f^*N$ has positive degree. As
\[
 L^2=p^2n+2pn=p(p+2)n>0,
\]
Nakai--Moishezon proves ampleness.

Apply Lemma~\ref{lem:eigen} with $a=p$. Only indices zero and one have nonnegative degree on a $\pi$-fibre; both have degree zero. This gives \eqref{eq:pushL}. Sections in the index-zero summand generate away from $E$. At $E$, the index-one generator is $u/s$, a unit times $u^{-p}$, and thus generates $\OO_S(pE)$ there. Its coefficient bundle is $N^{p+1}$. Both coefficient bundles are generated, so $L$ is generated everywhere.

For $r\ge1$, write
\[
 m_i(r)=\left\lceil\frac{p(r+i)}{p+1}\right\rceil.
\]
Formula~\eqref{eq:eigen} gives
\begin{equation}\label{eq:negative}
 \psi_*L^{-r}=\bigoplus_{i=0}^{p}\OO_P(-m_i(r))\otimes\pi^*N^{pi-r}.
\end{equation}
All $m_i(r)$ are positive, hence $f_*L^{-r}=0$. Relative duality on $\PP_C(V)$ gives, for $m\ge2$,
\begin{equation}\label{eq:duality}
 R^1\pi_*\OO_P(-m)=(\Sym^{m-2}V)^\vee\otimes T^{-1}.
\end{equation}
We always dualize after taking the symmetric power. For $r=1$ and $i=p$, the resulting summand of $R^1f_*L^{-1}$ is
\[
 (\Sym^{p-2}V)^\vee\otimes T^{p-2}.
\]
The quotient $V\to T$ gives an injection of $\OO_C$ into this bundle. Leray and $f_*L^{-1}=0$ therefore show that $H^1(S,L^{-1})\ne0$.

It remains to prove the claimed very ample power. Put $q=\ell^2$. Then
\[
 L^q=\psi^*F,\qquad F=\OO_P(p\ell)\otimes\pi^*N^q.
\]
For $a\ge1$, the bundle $\OO_P(a)\otimes\pi^*N^b$ is very ample if $bn\ge2g+1$. Indeed $\Sym^aV\otimes N^b$ has a filtration with quotients $N^{b+\ell j}$, $0\le j\le a$. After subtracting any subscheme of $C$ of length two, each quotient has degree greater than $2g-2$ and has no $H^1$. The vector bundle therefore supplies all values and first jets on the base. On a fibre the complete linear system of $\OO(a)$ separates points and tangent vectors. These two facts separate any pair of points and any tangent vector on $P$.

Here $2g-2=p\ell n$ and $qn-(2g+1)=\ell n-3>0$. Thus $F$ is very ample and $F\otimes M=\OO_P(p\ell-1)\otimes\pi^*N^{q+p}$ is generated. These imply that $\psi^*F$ is very ample: sections of $F$ recover the embedded base, and local ratios of sections from $F\otimes M$ recover a generator of the cyclic algebra. On each affine chart this gives a surjection onto the coordinate ring of the cover, and properness makes the resulting immersion closed.
\end{proof}

\begin{lemma}[Exact characteristic-two dimensions]\label{lem:dimensions}
For $p=2$,
\[
 h^1(S,L^{-r})=\begin{cases}1,&r=1,\\0,&r\ge2,\end{cases}
 \quad h^0(S,L)=6,\quad h^0(S,L^9)=556.
\]
\end{lemma}
\begin{proof}
Lemma~\ref{lem:curve} gives $g=7$ and $N=\OO_C(2P_\infty)$. The affine ring $k[x,y]/(y^2+y-x^{15})$ is free over $k[x]$ on $1,y$. Its monomials have distinct pole orders $2a$ and $2a+15$. Thus
\[
 h^0(N)=2,\qquad h^0(N^3)=4,
\]
and \eqref{eq:pushL} gives $h^0(L)=6$.

For negative powers, combine \eqref{eq:negative} and \eqref{eq:duality}. The contributions to $H^1(S,L^{-r})$ are
\[
 H^0\bigl(C,(\Sym^{m_i(r)-2}V)^\vee\otimes N^{2i-r-3}\bigr)
 \quad(0\le i\le2),
\]
with terms of negative symmetric degree omitted. At $r=1$, $m_i(r)$ equals $1,2,2$, so the three contributions are $0$, $H^0(N^{-2})$, and $H^0(\OO_C)$. The dimension is one. At $r\ge2$, each additional exponent $2i-r-3$ is at most $1-r<0$. The dual symmetric power has a filtration with line quotients $N^{-3j}$; all line quotients after twisting have negative degree. No contribution has a section.

For $L^9$, the three pushforward bundles on $C$ are
\[
 \Sym^6V\otimes N^9,\qquad
 \Sym^5V\otimes N^{11},\qquad
 \Sym^4V\otimes N^{13}.
\]
Their filtration quotients have degrees $2(9+3a)$, $2(11+3a)$, and $2(13+3a)$, all greater than $12=2g-2$. Riemann--Roch and the vanishing of their first cohomology yield
\begin{align*}
 h^0(L^9)
 &=\sum_{a=0}^{6}(12+6a)+\sum_{a=0}^{5}(16+6a)
                      +\sum_{a=0}^{4}(20+6a)\\
 &=210+186+160=556.
\end{align*}
These are finite exact dimensions, not numerical approximations.
\end{proof}

\section{Simultaneous roots and smooth hyperplane sections}\label{sec:roots}

The following elementary construction converts the generated polarization into a very ample one without losing its nonvanishing.

\begin{lemma}[Very ampleness by simultaneous roots]\label{lem:roots}
Let $S$ be a smooth projective connected surface over an algebraically closed field of characteristic $p$, and let $L$ be ample and generated. If $q>1$ is prime to $p$ and $L^q$ is very ample, choose a general basis $b_1,\ldots,b_r$ of $H^0(S,L^q)$ and form the simultaneous root cover $\rho:X\to S$ with $z_i^q=b_i$. Then $X$ is smooth, projective, and connected, $\deg\rho=q^r$, and $A=\rho^*L$ is very ample. For every integer $d$,
\begin{equation}\label{eq:root-cohomology}
 H^1(X,A^{-d})=
 \bigoplus_{0\le a_1,\ldots,a_r<q}
 H^1\bigl(S,L^{-(d+a_1+\cdots+a_r)}\bigr).
\end{equation}
\end{lemma}
\begin{proof}
Bertini and avoidance of finitely many intersections allow a basis whose zero divisors are smooth, pairwise transverse, and have no triple intersection. These properties are open, and intersect the open set of bases in the space of $r$-tuples of sections. Each divisor is connected because it is ample; being smooth, it is irreducible.

The defining algebra is locally free of rank $q^r$. At a point on one or two branch divisors their equations form part of a regular system of parameters. Adjoining roots replaces those parameters by $q$th roots; the other equations define étale extensions because their right sides are units and $q$ is invertible. Hence $X$ is smooth. A finite morphism to a projective scheme is projective here, since the pullback of an ample bundle is ample. Moreover
\[
 H^0(X,\OO_X)=\bigoplus_{0\le a_i<q}H^0(S,L^{-\sum a_i})=k.
\]
Thus $X$ is connected. A smooth connected scheme over a field has only one irreducible component, so it is integral.

Let $s_0,\ldots,s_v$ generate $L$. The sections $\rho^*s_j$ and $z_i$ define a map to projective space. On the chart $s_j\ne0$, the functions $b_i/s_j^q$ generate the coordinate ring of $S_{s_j}$, because $L^q$ embeds $S$. On the cover they equal $(z_i/s_j)^q$, and the ratios $z_i/s_j$ generate the remaining finite algebra. Thus the projective map is a closed immersion on these charts. They cover $X$, and properness gives a global closed immersion. The hyperplane bundle is $A$. Finally the splitting of $\rho_*\OO_X$ and the projection formula give \eqref{eq:root-cohomology}.
\end{proof}

\begin{proof}[Proof of Theorems~\ref{thm:main}(1)--(2) and the surface assertion of Theorem~\ref{thm:all-primes}]
Use Proposition~\ref{prop:polarization} with $q=(p+1)^2$ and Lemma~\ref{lem:roots}. The summand with all $a_i=0$ preserves the nonzero $H^1(S,L^{-1})$. In characteristic two, Lemma~\ref{lem:dimensions} shows that this is the only nonzero summand for $d=1$, and none survive for $d\ge2$. There are six sections of $L$ and 556 root sections, so the displayed immersion has target $\PP^{561}$ and degree $9^{556}$ over $S$.

Every smooth divisor $Y_d\in|A^d|$ is connected because it is ample. Since $H^0(X,\OO_X)=H^0(Y_d,\OO_{Y_d})=k$, its divisor sequence identifies the restriction kernel in $H^1(\OO)$ with $H^1(X,A^{-d})$. This proves all the assertions. Smooth sections exist by Bertini.
\end{proof}

All choices can be made over the algebraic closure of $\mathbf F_p$. Alternatively, a generic invertible matrix applied to any basis, followed by a generic hyperplane, specifies the construction over the algebraic closure of a rational function field. The assertion is about every smooth member once the embedding has been fixed.

\section{Differentials and the numerical constant}\label{sec:differentials}

We next bound the number of global one-forms, and prove the twisted vanishing that will control higher-degree integral restriction.

\begin{lemma}[Bounds on one-forms]\label{lem:differentials}
For the characteristic-two surface $X$,
\[
 h^0(X,\Omega_X^1)\le591\cdot9^{556},\qquad
 H^0(X,\Omega_X^1\otimes A^{-r})=0\quad(r\ge3).
\]
For the surfaces in characteristic $p$ constructed above,
\begin{equation}\label{eq:uniform-vanishing}
 H^1(X,A^{-r})=0\ (r\ge p^2-p),\qquad
 H^0(X,\Omega_X^1\otimes A^{-r})=0\ (r\ge p+1).
\end{equation}
\end{lemma}
\begin{proof}
For a tame root cover with simple normal-crossing branch divisor $D$, there is an injection
\[
 \Omega^1_{\text{cover}}\longrightarrow\rho^*\Omega^1_{\text{base}}(\log D).
\]
Locally, if $u^q=s$, then $ds/s=q\,du/u$ and $q$ is invertible. The injection sends $du$ to $(u/q)\,ds/s$; at a crossing use both branch parameters. These local formulas glue through the identification on the complement of the branch.

Apply this first to $\psi:S\to P$ in characteristic two. The relative cotangent sequence gives $h^0(P,\Omega_P^1)=g=7$ and
$H^0(P,\Omega_P^1\otimes M^i)=0$ for $i=1,2$: the two relative graded terms have fibre degrees $-i$ and $-2-i$. The residue sequence for $E_0+G$ bounds the untwisted logarithmic contribution by $7+2=9$. For $i=1,2$, the residue at $E_0$ is $N^{-i}$, which has no section. At $G$ it is $(\pi|_G)^*N^{2i}$, of degree $8i$. A line bundle of nonnegative degree $e$ on a smooth integral curve has at most $e+1$ sections, by successive subtraction of points. Hence
\[
 h^0(S,\Omega_S^1)\le9+9+17=35.
\]

The second cover has 556 smooth irreducible branch components. Its residue sequence therefore gives
\[
 h^0\bigl(S,\Omega_S^1(\log\textstyle\sum D_i)\bigr)\le35+556=591.
\]
The projection formula bounds $h^0(X,\Omega_X^1)$ by the sum, over $9^{556}$ indices, of the same logarithmic sheaf twisted by $L^{-\sum a_i}$. Multiplication by a nonzero section of $L^{\sum a_i}$ injects each such twist into the untwisted locally free sheaf. This proves the first inequality.

For the first vanishing in \eqref{eq:uniform-vanishing}, the summands in \eqref{eq:negative} and relative duality have filtration quotients with largest exponent
\[
 pi-r-(p+1)\le p^2-r-p-1<0
\]
when $r\ge p^2-p$. Their global sections vanish. Projection formula passes the vanishing from $S$ to $X$.

For one-forms, the exact differential defining $G$ gives an exact sequence
\begin{equation}\label{eq:ruled-log}
 0\longrightarrow\OO_P(p)\longrightarrow
 \Omega_P^1(\log(E_0+G))\longrightarrow
 \OO_P(-1)\otimes\pi^*T\longrightarrow0.
\end{equation}
Indeed the left term is $\pi^*\Omega_C^1(G)$ by \eqref{eq:branch} and $\Omega_C^1=T^p$. Locally $G$ has equation $z^p-f_i$ with $df_i\ne0$: its logarithmic pole enlarges the base differential line by $G$, because its relative derivative is zero. The disjoint section $E_0$ enlarges the relative line by $E_0$. This also proves exactness at the branch points.

Tensor the logarithmic inclusion for $\psi$ with $L^{-r}$ and use \eqref{eq:negative}. The right graded terms in \eqref{eq:ruled-log} have negative fibre degree. The left ones are
\[
 \OO_P(p-m_i(r))\otimes\pi^*N^{pi-r}.
\]
At $r=p+1$, only $i=0$ has nonnegative fibre degree, and that term is $\pi^*N^{-(p+1)}$, without sections. At larger $r$ all fibre degrees are negative. Thus the claimed vanishing holds on $S$.

For the multiroot cover, the additional residue terms after a negative twist are negative powers of an ample bundle on each $D_i$ and have no sections. Consequently the relevant logarithmic global sections equal the ordinary global sections, already zero on $S$. Projection formula proves the same vanishing on $X$.
\end{proof}

\section{From coherent nonvanishing to crystalline torsion}\label{sec:torsion}

This section supplies the formal cohomological mechanism of the integral counterexample. No degeneration of the Hodge spectral sequence is required.

\begin{lemma}[First-degree dimensions]\label{lem:torsion-bounds}
For a smooth projective geometrically connected variety $V$ over a perfect field of characteristic $p$, put
\[
 \begin{aligned}
 q(V)&=h^1(V,\OO_V),& u(V)&=h^0(V,\Omega_V^1),\\
 b(V)&=\rank_W C^1(V),& t(V)&=\dim_k C^2(V)[p].
 \end{aligned}
\]
Then $C^1(V)$ is free and
\begin{equation}\label{eq:torsion-bounds}
 q(V)-b(V)\le t(V)\le q(V)+u(V)-b(V).
\end{equation}
\end{lemma}
\begin{proof}
Derived crystalline base change gives the natural exact sequences
\begin{equation}\label{eq:uc}
 0\longrightarrow C^i(V)/p\longrightarrow H^i_{\dR}(V/k)
 \longrightarrow C^{i+1}(V)[p]\longrightarrow0.
\end{equation}
In degree zero, the first map is $k\xrightarrow{\sim}k$. Hence $C^1(V)[p]=0$, and finiteness over $W$ implies that $C^1(V)$ is free. In degree one,
\[
 h^1_{\dR}(V)=b(V)+t(V).
\]
The Hodge spectral sequence bounds this dimension above by $q(V)+u(V)$. The conjugate spectral sequence, through Cartier's isomorphism, has a natural edge injection
\[
 H^1(V^{(p)},\OO_{V^{(p)}})\longrightarrow H^1_{\dR}(V/k).
\]
Its term of bidegree $(1,0)$ admits no differential. Its dimension is $q(V)$ because $k$ is perfect. This gives the lower bound. The exact source for the Cartier and conjugate statements is \citet[0, Theorem 2.1.9, p.~515; II, \S6.14, p.~649]{Illusie}.
\end{proof}

\begin{proposition}[Product torsion growth]\label{prop:product-growth}
Let $X$ be a smooth projective connected surface and $A$ a very ample line bundle with $a=h^1(X,A^{-1})>0$. Let $m\ge1$, $Z=X\times\PP^1$, and $B=A\boxtimes\OO(m)$. For every smooth $Y\in|B|$,
\begin{equation}\label{eq:growth}
 t(Y)-t(Z)\ge a(m-1)-u(X).
\end{equation}
If the right side is positive, $T^2(Z)\to T^2(Y)$ is not surjective and $\coker\{C^2(Z)\to C^2(Y)\}$ has nonzero torsion.
\end{proposition}
\begin{proof}
Künneth gives $H^1(Z,B^{-1})=0$ and
\[
 H^2(Z,B^{-1})=H^1(X,A^{-1})\otimes H^1(\PP^1,\OO(-m)),
\]
of dimension $a(m-1)$. Multiplication by the equation of $Y$ gives the zero map from this group to $H^2(Z,\OO_Z)$. To see this, write the equation as a sum of tensors of sections of $A$ and $\OO(m)$. Each acts on the second Künneth factor through $H^1(\PP^1,\OO)=0$. The divisor sequence therefore gives
\[
 0\to H^1(X,\OO_X)\to H^1(Y,\OO_Y)
 \to H^1(X,A^{-1})\otimes H^1(\PP^1,\OO(-m))\to0.
\]
Thus $q(Y)=q(X)+a(m-1)$ and $u(Z)=u(X)$. Rational crystalline weak Lefschetz gives $b(Y)=b(Z)=b(X)$ and injectivity on rational second cohomology. Applying \eqref{eq:torsion-bounds} to $Y$ and $Z$ proves \eqref{eq:growth}.

For a finite torsion $W$-module, the dimension of its subgroup killed by $p$ equals its minimum number of generators. A quotient cannot require more generators than its source, so the strict inequality prevents surjectivity on torsion. Rational injectivity also gives
\[
 \operatorname{im}(C^2(Z)\to C^2(Y))\cap T^2(Y)
       =\operatorname{im}(T^2(Z)\to T^2(Y)).
\]
Consequently the nonzero torsion quotient of $T^2(Y)$ injects into the full cokernel. This proves the assertion.
\end{proof}

\begin{proof}[The degree-one assertions of Theorems~\ref{thm:main} and~\ref{thm:all-primes}]
In characteristic two, $a=1$ and Lemma~\ref{lem:differentials} gives $u(X)\le591\cdot9^{556}$. The prescribed $m$ makes the right side of \eqref{eq:growth} at least one. This proves torsion growth and the failure of saturation. In every characteristic the surface construction gives $a\ge1$, and choosing $m=u(X)+2$ suffices. The next section proves integral injectivity for these product sections and the exact higher-degree assertions.
\end{proof}

For completeness, the choice of $m$ can be made entirely numerical in every prime. With $n$, $g$, $\ell$ as in Lemma~\ref{lem:curve}, put $q=\ell^2$ and
\begin{align}
 B_p&=g+p+2+\frac{p^3(p+1)n}{2},\label{eq:Bp}\\
 a_i&=p\ell-i,\qquad
 r=\sum_{i=0}^{p}\left[(a_i+1)((q+pi)n+1-g)
                  +\frac{\ell n\,a_i(a_i+1)}2\right].\label{eq:r}
\end{align}
Then $r=h^0(S,L^q)$. Indeed its filtered line quotients are $N^{q+pi+\ell j}$ for $0\le j\le a_i$, all of degree greater than $2g-2$. Riemann--Roch gives \eqref{eq:r}. The same residue computation as above gives
\[
 h^0(S,\Omega_S^1)\le g+2+\sum_{i=1}^{p}(p^2in+1)=B_p.
\]
The multiroot cover consequently has $u(X)\le q^r(B_p+r)$. Thus
\[
 m=q^r(B_p+r)+2
\]
is an explicit sufficient choice. At $p=2$, $B_2=35$, $r=556$, and this is exactly $591\cdot9^{556}+2$.

\section{Integral injectivity and the exact degree threshold}\label{sec:degrees}

We now prove saturation in degree at least two for the fixed example. The relative de Rham differential is essential: vanishing of coherent kernels alone does not prove saturation of an integral crystalline map.

\begin{proposition}[A product criterion]\label{prop:criterion}
Let $X$ be a smooth projective connected surface, $A$ very ample, $m\ge1$, $Z=X\times\PP^1$, and $B=A\boxtimes\OO(m)$. For every smooth $Y\in|B^d|$, $d\ge1$, the map $C^2(Z)\to C^2(Y)$ is injective. If
\begin{equation}\label{eq:criterion}
 H^1(X,A^{-jd})=0\ (j=1,2,3),\qquad
 H^0(X,\Omega_X^1\otimes A^{-2d})=0,
\end{equation}
then its cokernel is torsion-free and it identifies the full torsion subgroups.
\end{proposition}
\begin{proof}
For the inclusion $i:Y\hookrightarrow Z$ let $K^\bullet$ be the kernel of $\Omega_Z^\bullet\to i_*\Omega_Y^\bullet$. In local coordinates with equation $s=0$ it is generated in degree $j$ by $s\Omega_Z^j$ and $ds\wedge\Omega_Z^{j-1}$. Consequently
\[
 K^j=\Omega_Z^j(\log Y)(-Y),\qquad
 0\to\Omega_Z^j(-Y)\to K^j\to\Omega_Y^{j-1}(-Y)\to0.
\]
In particular $K^0=\OO_Z(-Y)$. Künneth gives, for every $r>0$,
\begin{align*}
 H^1(Z,B^{-r})&=0,\\
 H^2(Z,B^{-r})&=H^1(X,A^{-r})\otimes H^1(\PP^1,\OO(-mr)),\\
 H^0(Z,\Omega_Z^j\otimes B^{-r})&=0\quad(j=1,2),\\
 H^1(Z,\Omega_Z^1\otimes B^{-r})
 &=H^0(X,\Omega_X^1\otimes A^{-r})\otimes H^1(\PP^1,\OO(-mr)).
\end{align*}
Since $\OO_Y(-Y)$ is antiample it has no sections. The residue sequence implies $H^0(K^1)=0$. Together with $H^1(K^0)=0$, this gives $\HH^1(Z,K^\bullet)=0$, hence an injection $H^1_{\dR}(Z)\to H^1_{\dR}(Y)$.

The free modules $C^1(Z)$ and $C^1(Y)$ have the same rank by rational weak Lefschetz. Their reductions modulo $p$ inject naturally into first de Rham cohomology, so their restriction map is injective modulo $p$. A map of free $W$-modules of the same finite rank that is injective modulo $p$ is an isomorphism. The first universal-coefficient diagram \eqref{eq:uc} now identifies
\[
 \ker\{C^2(Z)[p]\to C^2(Y)[p]\}
       \simeq\ker\{H^1_{\dR}(Z)\to H^1_{\dR}(Y)\}=0.
\]
The full second-crystalline kernel is torsion by rational weak Lefschetz. Every nonzero finite $p$-primary module has a nonzero element killed by $p$, so this kernel is zero.

Assume \eqref{eq:criterion}. The divisor sequences at twists $-Y$ and $-2Y$ give
\[
 H^2(K^0)=0,\qquad
 H^1(Y,\OO_Y(-Y))=H^1(Y,\OO_Y(-2Y))=0.
\]
The last Künneth formula also gives $H^1(Z,\Omega_Z^1(-2Y))=0$. Restricting $\Omega_Z^1(-Y)$ and then using the twisted conormal sequence
\[
 0\to\OO_Y(-2Y)\to\Omega_Z^1|_Y(-Y)\to\Omega_Y^1(-Y)\to0
\]
gives $H^0(Y,\Omega_Y^1(-Y))=0$. Hence $H^0(K^2)=0$.

The only possible term of total degree two on the first hypercohomology page is therefore $H^1(K^1)$. The residue sequence identifies it with $H^1(Z,\Omega_Z^1(-Y))$. Its outgoing differential is injective. Indeed the composite
\[
 H^1(Z,\Omega_Z^1(-Y))\xrightarrow{\sim}H^1(K^1)
 \xrightarrow{d_1}H^1(K^2)\longrightarrow H^1(Y,\Omega_Y^1(-Y))
\]
is, up to sign, ordinary restriction. The local identity is
$d(s\eta)=ds\wedge\eta+s\,d\eta$; taking the residue leaves $\eta|_Y$ with its indicated twist. Restriction into $H^1(Y,\Omega_Z^1|_Y(-Y))$ is injective because $H^1(Z,\Omega_Z^1(-2Y))=0$. The next map is injective because $H^1(Y,\OO_Y(-2Y))=0$. Thus $d_1$ is injective, all total-degree-two terms vanish on the second page, and $\HH^2(Z,K^\bullet)=0$.

It follows that first de Rham restriction is an isomorphism and second de Rham restriction is injective. Set $M=C^2(Z)$ and $N=C^2(Y)$. We already know $M\hookrightarrow N$. The first universal-coefficient diagram gives $M[p]\xrightarrow{\sim}N[p]$, while the second gives $M/pM\hookrightarrow N/pN$. For $Q=N/M$, the exact Tor sequence contains
\[
 M[p]\longrightarrow N[p]\longrightarrow Q[p]
             \longrightarrow M/pM\longrightarrow N/pN.
\]
The outside maps are respectively surjective and injective, so $Q[p]=0$. Therefore $Q$ is torsion-free and $M$ and $N$ have the same torsion subgroup under restriction.
\end{proof}

\begin{proof}[Completion of Theorems~\ref{thm:main} and~\ref{thm:all-primes}]
The product criterion proves integral injectivity in every positive degree. In characteristic two, if $d\ge2$, all three negative first-cohomology groups in \eqref{eq:criterion} vanish by Lemma~\ref{lem:dimensions} and the root decomposition. The one-form group vanishes by Lemma~\ref{lem:differentials}, since $2d\ge4$. Hence the cokernel is torsion-free and the torsion subgroups agree. In degree one, Proposition~\ref{prop:product-growth} shows that both conclusions fail. This proves the two equivalences.

For general $p$, both vanishings in \eqref{eq:criterion} hold whenever $d\ge p^2-p$, by \eqref{eq:uniform-vanishing}; indeed $2(p^2-p)\ge p+1$. Thus the same product construction is saturated in every such degree and nonsaturated in degree one.
\end{proof}

\section{A threefold with an integral restriction kernel}\label{sec:injectivity}

We now construct a second threefold. Its role differs from that of the product: integral restriction itself fails in the first two degrees. All varieties in this section are over an algebraically closed field of characteristic two. A \emph{Tango datum} on a smooth projective variety $B$ means an effective ample divisor $D$ and a rational function $f$ such that $df$, divided locally by an equation of $2D$, is a regular nowhere-zero one-form. The stronger nowhere-zero condition permits iteration while preserving smoothness. We use the iteration of \citet[Proposition 2.3, p.~8 of the preprint]{Mukai}, but give the local construction in the quotient convention used throughout this paper.

\begin{lemma}[One iteration]\label{lem:iteration}
Suppose $(B,D,f)$ is a Tango datum and $D=qN$, where $q>2$ is odd and $N$ is effective. There are a smooth projective variety $\widetilde B$, a flat morphism $g:\widetilde B\to B$ with geometrically integral curve fibres, and a section $E$, with a Tango datum
\[
 \widetilde D=(q-1)E+g^*N.
\]
Let $0\to\OO_B\to V\to\OO_B(D)\to0$ be the nonzero Frobenius-killed extension defined by $df$, and set $P=\PP_B(V)$. The map $g$ factors through a finite flat map $\psi:\widetilde B\to P$ of degree $q$. For every integer $a$,
\begin{equation}\label{eq:iteration}
 \psi_*\OO_{\widetilde B}(aE)=
 \bigoplus_{i=0}^{q-1}\OO_P\!\left(\left\lfloor\frac{a-2i}{q}\right\rfloor\right)
             \otimes\pi^*\OO_B(2iN).
\end{equation}
For each prime $\ell\ne2$, pullback induces
$H^1_{\et}(B,\mathbf Q_\ell)\simeq H^1_{\et}(\widetilde B,\mathbf Q_\ell)$.
\end{lemma}
\begin{proof}
The Cartier exact-differential sequence, twisted by $\OO_B(-D)$, sends the normalized differential to a class in $H^1(B,\OO_B(-D))$ killed by Frobenius. It is nonzero because the preceding group $H^0(B,\OO_B(-2D))$ vanishes. This is the same construction as for the curve in Section~\ref{sec:surface}; in higher dimension the exact sequence is given in \citet[Theorem 1 and Lemma 1.4, pp.~1, 4--5 of the preprint]{Mukai}.

Let $b_i$ be local equations for $N$, and put $a_i=b_i^q$. Cartier gives, after restricting the open sets, regular functions $f_i$ and rational functions $h_i$ such that
\[
 f/a_i^2=f_i+h_i^2,\qquad df_i=a_i^{-2}df.
\]
The $df_i$ are nowhere zero. On overlaps write $t_{ij}=a_i/a_j$ and $c_{ij}=t_{ij}h_i-h_j$. Then
$f_j=t_{ij}^2f_i+c_{ij}^2$. The functions $c_{ij}$ are regular, since their squares are regular and the rings are normal. Glue the affine charts of $P$ by
$S_j=t_{ij}S_i+c_{ij}$. Their infinity section $E_0$ has normal bundle $\OO_B(D)$ and divisor bundle $\OO_P(1)$. The equations $S_i^2=f_i$ define a smooth divisor $G$ disjoint from $E_0$, with
\[
 \OO_P(G)=\OO_P(2)\otimes\pi^*\OO_B(-2D).
\]
Smoothness follows from the nowhere-zero $df_i$.

Set $M=\OO_P(-1)\otimes\pi^*\OO_B(2N)$. The relation
$M^{-q}=\OO_P(G+(q-2)E_0)$ defines a cyclic cover; take its normalization. Near $G$ the branch multiplicity is one. Near $E_0$, after an étale extension extracting a root of a unit, its equation is $w^q=s^{q-2}$, whose normalization has $s=u^q$, $w=u^{q-2}$. These are regular local models because $\gcd(q,q-2)=1$. They prove smoothness and $\psi^*E_0=qE$. The equation is Eisenstein at $G$, so the cover is integral. In eigendegree $i$, normalization and twisting by $aE$ allow the pole $\lfloor((q-2)i+a)/q\rfloor E_0$. Thus its eigensheaf is
\[
 M^i\bigl(\lfloor((q-2)i+a)/q\rfloor E_0\bigr),
\]
which is precisely \eqref{eq:iteration}. The algebra is locally free, so $\psi$ and $g$ are flat.

On the affine chart the normalized cover has equation
$T_i^q=S_i^2-f_i$. Its transition is $T_j=(b_i/b_j)^2T_i$, and hence $\widetilde f=b_i^2T_i$ is a global rational function. Since
\[
 d\widetilde f=b_i^2dT_i,\qquad qT_i^{q-1}dT_i=-df_i,
\]
the normalized differential is nonzero off the branch. At $G$, one can use $f_i$ as a base coordinate and solve the equation for it; $T_i$ is then a smooth parameter and $dT_i$ is nonzero. At $E$, the function $T_i$ has pole order two, while $df_i$ pulls back to a nowhere-zero base differential. The displayed identity therefore gives exact zero order $2(q-1)$ for $dT_i$. This proves the asserted Tango condition for $\widetilde D$.

A geometric fibre has affine equation $T^q=(S-c)^2$ and one point at infinity. Its normalization is $\PP^1$, and that normalization is a universal homeomorphism. The fibre is integral, and $E$ has degree one on it. Moreover
$(\OO(E)|_E)^q\simeq\OO_B(qN)$, so $E|_E$ is numerically ample. The effective divisor $E$ is nef: curves outside $E$ intersect it nonnegatively, and its restriction to $E$ is ample. For a $d$-dimensional subvariety $Z$ of $\widetilde B$, the image in $B$ has dimension $d$ or $d-1$. In the first case $(g^*N)^d\cdot Z>0$; in the second, $E(g^*N)^{d-1}\cdot Z>0$. All other terms in $\widetilde D^d\cdot Z$ are nonnegative. Nakai--Moishezon proves ampleness.

Finally each geometric fibre has the étale cohomology of $\PP^1$. Proper base change gives $g_*\mathbf Q_\ell=\mathbf Q_\ell$ and $R^1g_*\mathbf Q_\ell=0$. The degree-one Leray sequence proves the last assertion. These uses of proper base change may be applied with $\mathbf Z/\ell^n$ first and then passed to $\mathbf Q_\ell$ \citep[Exposé XII, Theorem 5.1 and Corollary 5.2, pp.~88--89]{SGA4}.
\end{proof}

\begin{definition}[The second threefold]\label{def:threefold}
Start from the smooth completion of $C:y^2+y=x^{843}$. Let $P_\infty$ be its point at infinity and $N_0=\OO_C(2P_\infty)$. The function $x$ has
$(dx)=840P_\infty$, so $D_1=420P_\infty=210N_0$ is a Tango divisor. Apply Lemma~\ref{lem:iteration} with $q=21$ and $N=10N_0$, obtaining $f:S\to C$ and $E_1$. Put
\[
 B=\OO_S(2E_1)\otimes f^*N_0.
\]
The new Tango divisor is $10B$. Apply the lemma again with $q=5$ and $N=2B$, obtaining $g:R\to S$ and $E_2$, and put
\[
 L=\OO_R(2E_2)\otimes g^*B.
\]
Its new Tango divisor is $2L$. Choose 843 general sections of $L^{445}$, with the first eight defining a closed immersion in $\PP^7$, and take their simultaneous 445th-root cover $\rho:U\to R$. Put $H=\rho^*L$.
\end{definition}

We prove all the existence and positivity assertions implicit in this definition before using it. Letters $U,H$ distinguish this threefold from the product $Z,B$ of Theorem~\ref{thm:main}; within this section $B$ denotes the surface polarization in Definition~\ref{def:threefold}.

\begin{proposition}[Positivity and first Betti number]\label{prop:threefold-positive}
The curve in Definition~\ref{def:threefold} has genus $421$. The line bundles $B$ and $L$ are ample and generated, $B^{441}$ and $L^{445}$ are very ample, and
\[
 h^0(S,B)=24,\qquad h^0(S,B^5)=336.
\]
The root cover $U$ is smooth, projective, and connected, $H$ is very ample, and
\[
 \deg\rho=445^{843},\qquad
 b_1(U):=\dim H^1_{\et}(U,\mathbf Q_3)=842.
\]
\end{proposition}
\begin{proof}
The calculation in Lemma~\ref{lem:curve}, with $h=422$, gives $g(C)=421$ and the stated divisor of $dx$. The sections $1,x$ generate $N_0$. Formula~\eqref{eq:iteration} gives
\[
 f_*B=N_0\oplus N_0^{21},\qquad g_*L=B\oplus B^5.
\]
In both steps the index-zero sections generate away from the ramification section; the index-one local generator $w/s=u^{-2}$ is nonvanishing in the frame of $\OO(2E)$ on that section. The coefficient bundles are generated. Hence $B$ and $L$ are generated. They are ample because $10B$ and $2L$ are the ample divisors of the iteration.

The extension on the first projective bundle is
$0\to\OO_C\to V_1\to N_0^{210}\to0$. For $a\ge1$ and $b\ge441$, the bundle $\Sym^aV_1\otimes N_0^b$ has a filtration with quotients $N_0^{b+210j}$. Each has degree at least $882>2g(C)+1$. Serre duality shows that its $H^1$ vanishes after subtracting every length-two subscheme. Sections therefore provide all base values and first jets, while $\OO(a)$ provides the complete fibre system. Thus $\OO_{\PP(V_1)}(a)\otimes N_0^b$ is very ample.

Now $B^{441}=\psi_1^*F_1$, where $F_1=\OO(42)\otimes N_0^{441}$. The cover algebra has summands
\[
 K_i=\OO\!\left(-\left\lceil2i/21\right\rceil\right)\otimes N_0^{20i},
 \quad 0\le i\le20.
\]
Every $F_1K_i$ is generated: its relative degree is at least 40 and its base exponent at least 441. Sections from $F_1$ recover the embedded projective bundle, and generators from all $F_1K_i$ recover its finite cover algebra on each affine chart. This proves very ampleness of $B^{441}$. More generally, for $t\ge441$ the eigensheaves of $B^t$ push to
\[
 \Sym^{\lfloor2(t-i)/21\rfloor}V_1\otimes N_0^{t+20i}.
\]
They have nonnegative relative degree and filtrations with line quotients of degree at least $882>2g-2$. Consequently $H^1(S,B^t)=0$. Since $B^t$ is very ample, its evaluation on every length-two subscheme $Z$ is surjective, and
\begin{equation}\label{eq:B-jets}
 H^1(S,B^t\otimes I_Z)=0\quad(t\ge441,\ \operatorname{length}Z\le2).
\end{equation}

The second extension is $0\to\OO_S\to V_2\to B^{10}\to0$. By \eqref{eq:B-jets}, the filtered bundle $\Sym^aV_2\otimes B^b$ evaluates surjectively on every length-two subscheme of $S$ whenever $a\ge1$, $b\ge441$. Fibre sections separate the remaining vertical directions. Hence $\OO_{\PP(V_2)}(a)\otimes B^b$ is very ample in this range. Since
\[
 L^{445}=\psi_2^*\bigl(\OO(178)\otimes B^{445}\bigr),
\]
and the normalized algebra summands are
$\OO(-\lceil2j/5\rceil)\otimes B^{4j}$, all their twists by $\OO(178)\otimes B^{445}$ are generated. The same finite-algebra argument proves that $L^{445}$ is very ample.

For $0\le b\le105$, the pole semigroup of $C$ gives $h^0(N_0^b)=b+1$: only $1,x,\ldots,x^b$ have pole order at most $2b<843$. Thus
\[
 h^0(B)=2+22=24,\qquad
 h^0(B^5)=\sum_{i=0}^5(6+20i)=336.
\]
There are at least 179 independent sections of $L^{445}$, since its index-zero pushforward has 179 positive line quotients with vanishing $H^1$. A general projection of the resulting threefold embedding to $\PP^7$ remains a closed immersion: the closed secant and tangent locus has dimension at most seven, and a general centre of codimension eight misses it. Thus eight sections can be chosen to embed $R$. Independently, 843 general divisors are smooth and meet with simple normal crossings, with no fourfold intersection. Both requirements are nonempty open conditions on an irreducible parameter space, so they hold together.

The root-cover local coordinates prove smoothness exactly as in Section~\ref{sec:roots}, now with at most three branch parameters. Its algebra is locally free of rank $445^{843}$. The decomposition into negative powers of the ample $L$ gives $H^0(U,\OO_U)=k$, so it is connected. A smooth connected variety is integral. The generated sections of $L$ together with the root sections embed $U$: on a chart $s\ne0$ of $L$, the functions $b_i/s^{445}$ from the first eight sections generate the base affine algebra after adjoining the section $s^{445}$ to their very ample system, and $(z_i/s)^{445}=b_i/s^{445}$. The ratios $z_i/s$ consequently generate the finite cover algebra. Properness finishes the closed-immersion assertion.

The two Tango steps preserve $H^1_{\et}$, so $b_1(R)=842$. To see that the root cover preserves this number, extend $\mathbf Q_3$ to a finite extension $K$ containing all 445th roots of unity and decompose $\rho_*K$ under $(\mu_{445})^{843}$. Its trivial character is $K_R$. For a nontrivial character $\chi$, let $J$ be the nonempty set of its nontrivial inertia indices. The corresponding sheaf is $j_!\mathcal L_\chi$, where $j:V=R\setminus\bigcup_{i\in J}D_i\hookrightarrow R$ and $\mathcal L_\chi$ is lisse of rank one. Locally this is the character decomposition of each tame root cover: a nontrivial character has zero stalk on its branch divisor, and a trivial one extends constantly. The open $V$ is smooth affine of dimension three because the removed divisor is ample. Affine vanishing and Poincaré duality give
\[
 H^1(R,j_!\mathcal L_\chi)=H_c^1(V,\mathcal L_\chi)
   =H^5(V,\mathcal L_\chi^\vee(3))^\vee=0.
\]
These are \citet[Exposé XIV, Theorem 3.1 and Corollary 3.2, p.~160; Exposé XVIII, Theorem 3.2.5 and (3.2.6.2), pp.~585--587]{SGA4}. Finite proper base change has no higher direct images for $\rho$. Thus only the trivial character contributes in degree one, proving $b_1(U)=842$.
\end{proof}

\subsection{The exact negative cohomology}

\begin{lemma}\label{lem:threefold-negative}
For the line bundles in Definition~\ref{def:threefold},
\[
 h^1(R,L^{-r})=\begin{cases}360&r=1,\\1&r=2,\\0&r\ge3,\end{cases}
 \qquad
 h^1(U,H^{-d})=\begin{cases}1203&d=1,\\1&d=2,\\0&d\ge3.\end{cases}
\]
\end{lemma}
\begin{proof}
Let $\alpha\in H^1(S,B^{-10})$ be the extension class of $V=V_2$. For $0\le e<10$,
\begin{equation}\label{eq:nonannihilation}
 H^0(S,V^\vee B^e)=0.
\end{equation}
Indeed this group is the kernel of multiplication by $\alpha$ from $H^0(B^e)$ to $H^1(B^{e-10})$. Under the Frobenius-kernel description, multiplying $\alpha$ by a nonzero section $s$ replaces its rational exact differential $\eta$ by $s^2\eta$. This remains nonzero. The connecting map for the ample divisor $(10-e)B$ is injective, since $H^0(B^{-2(10-e)})=0$. Thus $s\alpha\ne0$, proving \eqref{eq:nonannihilation}.

Put $m_j(r)=\lceil2(r+j)/5\rceil$ and $e_j(r)=4j-r-10$. Formula~\eqref{eq:iteration} and relative duality give
\begin{equation}\label{eq:threefold-negative}
 H^1(R,L^{-r})=
 \bigoplus_{m_j(r)\ge2}
 H^0\bigl(S,(\Sym^{m_j(r)-2}V)^\vee B^{e_j(r)}\bigr).
\end{equation}
All zeroth direct images vanish. The dual in this formula is taken after the symmetric power; commuting these operations is not valid in characteristic two.

For $r\ge7$, every $e_j(r)\le6-r<0$, and the dual symmetric power has only quotients $B^{-10a}$ in its filtration, so all sections vanish. At $r=3$, the only possible term is $H^0(V^\vee B^3)$, zero by \eqref{eq:nonannihilation}. At $r=4,5,6$, the only possible term is $H^0((\Sym^2V)^\vee B^{6-r})$. Frobenius splitting gives the decomposition
\begin{equation}\label{eq:sym-two}
 \Sym^2V\simeq V\oplus B^{20}.
\end{equation}
For clarity, if $u$ generates the trivial subline and $v_j=a_{ij}v_i+c_{ij}u$, the subbundle $uV$ is generated by $u^2,uv_i$ and is isomorphic to $V$. A splitting of $F^*V$ provides combinations $v_i^2+h_i u^2$ that glue as $B^{20}$. They complement $uV$. This proves \eqref{eq:sym-two}; its dual and \eqref{eq:nonannihilation} give the required vanishing.

At $r=2$, the $j=3$ term in \eqref{eq:threefold-negative} is $H^0(\OO_S)=k$, while the $j=4$ term is $H^0(V^\vee B^4)=0$. At $r=1$, the only nonzero terms are $H^0(B)$ and $H^0(B^5)$, of dimensions 24 and 336. This proves the first formula. Finally
\[
 H^1(U,H^{-d})=
 \bigoplus_{0\le a_i<445}H^1(R,L^{-(d+\sum a_i)}).
\]
There are 843 tuples with sum one. Hence the dimensions are $360+843=1203$, $1$, and $0$ in the asserted ranges.
\end{proof}

\subsection{Negative one-forms}

The second iteration supplies a nowhere-zero map $\vartheta:L^4\to\Omega_R^1$. Use the same letter for its pullback $H^4\to\Omega_U^1$; the latter is generically nonzero because the cover is separable.

\begin{lemma}\label{lem:threefold-forms}
For every $r\ge2$ and $d\ge2$,
\[
 H^0(R,\Omega_R^1L^{-r})=\vartheta H^0(R,L^{4-r}),\qquad
 H^0(U,\Omega_U^1H^{-d})=\vartheta H^0(U,H^{4-d}).
\]
The map $\vartheta:H^4\to\Omega_U^1$ does not vanish identically on any smooth divisor $Y\in|H^d|$, $d\ge1$.
\end{lemma}
\begin{proof}
The first iteration gives a subbundle $B^{20}\subset\Omega_S^1$ with line-bundle quotient $Q$. On $P=\PP_S(V)$, the logarithmic cotangent bundle for the disjoint branches $E_0,G$ has a filtration with quotients
\[
 \OO_P(2),\qquad \pi^*Q,\qquad \OO_P(-1)\otimes B^{10}.
\]
To check this locally, write $G$ as $S^2-f_i=0$, with $df_i$ spanning the distinguished base line. The horizontal logarithmic generator $df_i/(S^2-f_i)$ allows precisely the twist $B^{20}(G)=\OO_P(2)$; the remaining horizontal quotient is $Q$. The relative logarithmic quotient has a pole only at $E_0$, hence is $\Omega^1_{P/S}(E_0)=\OO_P(-1)B^{10}$. In the chart $z=1/S$, these generators become $df_i/(1-f_i z^2)$ and $dz/z$, giving the same filtration across infinity.

The tame normalized cover embeds $\Omega_R^1$ in the pullback of this logarithmic bundle. After twisting by $L^{-r}$, the last two pieces have relative degrees $-m_j(r)$ and $-1-m_j(r)$, and therefore no global sections. The first piece has eigensheaves
\[
 \OO_P(2-m_j(r))\otimes B^{4j-r}.
\]
For $r\ge6$ all relative degrees are negative. At $r=5$, the only degree-zero term is $B^{-5}$, which also has no sections. At $r=4$, the only nonzero possibility is $H^0(\OO_S)$. At $r=3$ the possibilities are $H^0(B)\oplus H^0(B^5)=H^0(R,L)$. At $r=2$ they are
\[
 H^0(VB^{-2})\oplus H^0(B^2)\oplus H^0(B^6)\oplus H^0(B^{10}).
\]
Here $H^0(VB^{-2})=0$: its extension sequence has connecting map $H^0(B^8)\to H^1(B^{-2})$, and \eqref{eq:nonannihilation} proves its injectivity. The other three terms equal $H^0(R,L^2)$ by \eqref{eq:iteration}. Thus in every case the logarithmic calculation bounds the left side by $h^0(L^{4-r})$. Multiplication by the nonzero $\vartheta$ injects that latter space into it, proving equality.

Let $D$ be the branch divisor of $\rho$. The residue sequence gives
\[
 H^0(R,\Omega_R^1(\log D)L^{-r})=H^0(R,\Omega_R^1L^{-r})
 \quad(r>0),
\]
since $L^{-r}|_{D_i}$ is antiample. The root-cover logarithmic inclusion and its eigensheaves therefore bound
\[
 h^0(U,\Omega_U^1H^{-d})
 \le\sum_{0\le a_i<445}h^0(R,L^{4-d-\sum a_i})=h^0(U,H^{4-d}).
\]
The pulled-back $\vartheta$ provides the opposite injection, and equality follows. If $\vartheta$ vanished identically along a divisor in $|H^d|$, its defining section of $\Omega_U^1H^{-4}$ would come from $H^0(\Omega_U^1H^{-4-d})$, which vanishes by the formula. This proves the last assertion without claiming that the pulled-back map is everywhere a subbundle.
\end{proof}

\subsection{The integral lattice in degree one}

\begin{lemma}\label{lem:first-lattice}
Let $h:U\to C$ be the composite morphism. Then $h_*\OO_U=\OO_C$, its geometric generic fibre is integral, and pullback gives
\[
 C^1(C)\xrightarrow{\sim}C^1(U).
\]
For every smooth $Y\in|H^d|$ with $d\ge2$, one also has $(h|_Y)_*\OO_Y=\OO_C$ and integral isomorphisms
\[
 C^1(C)\xrightarrow{\sim}C^1(Y),\qquad C^1(U)\xrightarrow{\sim}C^1(Y).
\]
\end{lemma}
\begin{proof}
Both Tango maps are flat with geometrically integral fibres. Their composite $R\to C$ has the same property: on a geometric fibre the second map is flat over the integral first fibre, so affine charts inject into their integral generic-fibre algebras, proving reducedness and irreducibility. Thus $k(R)/k(C)$ is regular. Let $U\to C'\to C$ be the Stein factorization. The curve $C'$ is smooth, and regularity implies that $k(R)$ and $k(C')$ are linearly disjoint over $k(C)$. Therefore
\[
 [k(C'):k(C)]\mid[k(U):k(R)]=445^{843}.
\]
In particular $C'\to C$ is separable. The map $U\to C'$ has geometrically connected fibres, so $H^1_{\et}(C',\mathbf Q_3)$ injects into $H^1_{\et}(U,\mathbf Q_3)$. Hence $g(C')\le421$. Riemann--Hurwitz gives
$g(C')-1\ge\deg(C'/C)(421-1)$, forcing degree one. This proves $h_*\OO_U=\OO_C$. The extension $k(U)/k(C)$ is separably generated, since $k(R)/k(C)$ is regular and the root cover is separable. It has no new algebraic constants, so it is regular and the geometric generic fibre is integral.

The two Tango maps have sections, and their pullbacks on reduced Picard varieties are split monomorphisms. Their first Betti numbers coincide, so these maps are isomorphisms. Thus $\operatorname{Pic}^0_{R,\mathrm{red}}\simeq\operatorname{Pic}^0_C$. For the finite locally free map $\rho$, the norm identity
\[
 \operatorname{Nm}_\rho\circ\rho^*=[445^{843}]
\]
shows that pullback on reduced Picard varieties is an isogeny with kernel killed by an odd integer. Here the dimensions agree by Proposition~\ref{prop:threefold-positive}. Its map on 2-divisible groups is therefore an isomorphism. The integral Dieudonné realization of the reduced Picard variety identifies these groups functorially with first crystalline cohomology \citep[II, Proposition 3.11 and Remark 3.11.2, pp.~618--619]{Illusie}. This proves the integral assertion for $U$.

Now let $d\ge2$. The divisor $Y$ dominates $C$: a divisor supported on finitely many fibres would miss another positive-dimensional fibre and could not be ample. The geometric generic fibre $U_{\overline\eta}$ is an integral projective Cohen--Macaulay surface, since it is a geometric fibre of a flat map from a smooth variety to a smooth curve. The restriction $Y_{\overline\eta}$ is an ample effective Cartier divisor. It is connected. Indeed, Serre duality and Serre vanishing give $H^1(\OO_{U_{\overline\eta}}(-nY))=0$ for large $n$, and the divisor sequence then gives $H^0(\OO_{nY_{\overline\eta}})=\overline{k(C)}$. A disconnected $Y_{\overline\eta}$ would yield a nontrivial idempotent on every $nY_{\overline\eta}$, contradicting this equality.

It follows that the finite Stein curve map $C_Y\to C$ is purely inseparable. We rule out a nontrivial such map by keeping track of the coefficient of the Tango differential. Write $h_R:R\to C$, and let $G_1'\subset S$ and $G_2'\subset R$ be the reduced inverse images of the two inseparable multisections. The effective divisor
\[
 D_0=210h_R^*P_\infty+5g^*G_1'+G_2'
\]
has line bundle $L$. Indeed,
\[
 \OO_S(G_1')=\OO_S(2E_1)\otimes f^*N_0^{-20},\qquad
 \OO_R(G_2')=\OO_R(2E_2)\otimes g^*B^{-4},
\]
so the first two summands of $D_0$ have line bundle $g^*B^5$, and the last changes it to $L$. Let $\sigma_0\in H^0(R,L)$ have divisor $D_0$. The local equations of the two covers give vanishing orders twenty and four for the successive pullbacks of the normalized base differential. Together with $(dx)=840P_\infty$, this gives
\[
 h_R^*dx=c\,\vartheta_R(\sigma_0^4),\qquad c\in k^\times,
\]
where $\vartheta_R:L^4\to\Omega_R^1$ is the second Tango subbundle. Both coefficients have divisor $4D_0$, so their ratio is a nonzero constant. Pullback gives the same identity on $U$.

The divisor $\rho^*D_0\in|H|$ cannot contain $Y\in|H^d|$ when $d\ge2$, because otherwise $H^{1-d}$ would have a nonzero section. Thus the coefficient of $\vartheta$ is nonzero at the generic point of $Y$. Moreover the image of $\vartheta$ in $\Omega_Y^1$ is nonzero: if it vanished, its restriction would factor through the conormal bundle $H^{-d}|_Y$, whereas $H^0(Y,H^{-d-4}|_Y)=0$; the absence of divisorial zeros of $\vartheta$ excludes a zero restriction. Therefore $(h|_Y)^*dx\ne0$, so $Y\to C$ is generically separable. The Stein map is both separable and purely inseparable and hence has degree one. This proves $C_Y=C$.

The equality $(h|_Y)_*\OO_Y=\OO_C$ makes $\operatorname{Pic}_C\to\operatorname{Pic}_Y$ a monomorphism of fppf Picard sheaves. To see this, the equality of structure-sheaf pushforwards persists after every base change over $k$, by flat base change. If a pulled-back line bundle is trivial, or comes from the parameter base, projection formula and pushforward recover the same assertion on $C$. After local fppf trivialization this proves the assertion for the sheaf quotient defining the Picard functor. Thus $\operatorname{Pic}^0_C\to\operatorname{Pic}^0_{Y,\mathrm{red}}$ is a closed immersion. Étale weak Lefschetz gives $b_1(Y)=b_1(U)=842$; equivalently, affine vanishing and duality on the smooth affine complement give the required degree-one isomorphism. The two abelian varieties both have dimension 421, so their closed immersion is an isomorphism. Their integral Dieudonné realizations give $C^1(C)\simeq C^1(Y)$. Factoring this map through $U$ proves the final assertion.
\end{proof}

\subsection{The exact injectivity range}

\begin{theorem}\label{thm:injectivity-exact}
For $(U,H)$ in Definition~\ref{def:threefold} and every smooth $Y_d\in|H^d|$, restriction
\[
 H^2_{\cris}(U/W)\longrightarrow H^2_{\cris}(Y_d/W)
\]
is injective if and only if $d\ge3$. In degree one its kernel killed by two has dimension at least 361, and in degree two this dimension is at least one. The coherent restriction kernels in $H^1(\OO)$ have dimensions $1203$, $1$, and $0$ for $d=1$, $d=2$, and $d\ge3$, respectively.
\end{theorem}
\begin{proof}
The coherent dimensions follow from Lemma~\ref{lem:threefold-negative}, the divisor sequence, and connectedness of each smooth ample section. The functorial conjugate edge injects this coherent kernel, after Frobenius twist, into the first de Rham restriction kernel. When $d=1$, its dimension is 1203. In the universal-coefficient sequence
\[
 0\to C^1(U)/2\to H^1_{\dR}(U/k)\to C^2(U)[2]\to0,
\]
the first term has dimension $842$. Naturality therefore supplies a subspace of dimension at least $1203-842=361$ in the kernel on $C^2[2]$.

For $d=2$, the coherent kernel is nonzero. In the natural diagram of the same universal-coefficient sequences for $U$ and $Y_2$, the left vertical map is an isomorphism by Lemma~\ref{lem:first-lattice}. The snake lemma identifies the first de Rham restriction kernel with the kernel on $C^2[2]$. The former contains the nonzero coherent kernel, proving failure of integral injectivity.

Let $d\ge3$, and first consider a global closed one-form $\omega$ on $U$ whose restriction to $Y=Y_d$ is zero. The conormal sequence has $H^0(Y,H^{-d}|_Y)=0$. Thus $\omega$ lies in the image of $H^0(\Omega_U^1H^{-d})$, multiplied by the defining section $s_Y$. For $d\ge5$ this space is zero by Lemma~\ref{lem:threefold-forms}. For $d=3,4$ write
\[
 \omega=s_Yt\vartheta,\qquad t\in H^0(U,H^{4-d}).
\]
If $t\ne0$, it cannot vanish on $Y$, since that would give a nonzero section of the antiample bundle $H^{4-2d}$. Locally write $s_Y=g$, $t=u$, and let $\vartheta_0$ represent $\vartheta$. Reduction of $d\omega=0$ modulo $g$ gives
$u\,dg\wedge\vartheta_0=0$. At the generic point of the integral smooth divisor $Y$, this implies that $\vartheta|_Y$ maps to zero in $\Omega_Y^1$. Since this latter bundle is torsion-free, the vanishing holds everywhere, and $\vartheta|_Y$ factors through the conormal line:
\[
 H^4|_Y\longrightarrow H^{-d}|_Y.
\]
There is no such nonzero homomorphism because $H^0(Y,H^{-d-4}|_Y)=0$. Thus $\vartheta$ itself would vanish along $Y$, contrary to Lemma~\ref{lem:threefold-forms}. Hence $\omega=0$.

For a smooth proper connected variety the Hodge edge identifies the kernel of $H^1_{\dR}\to H^1(\OO)$ with global closed one-forms; those forms inject into de Rham cohomology because the differential of every global function is zero. A de Rham restriction-kernel class has zero coherent edge, as coherent restriction is injective for $d\ge3$. It therefore comes from a closed global one-form restricting to zero, which the preceding paragraph proves is zero. Thus first de Rham restriction is injective.

The first crystalline lattices are isomorphic by Lemma~\ref{lem:first-lattice}. The universal-coefficient diagram now gives injectivity on $C^2[2]$. Rational crystalline weak Lefschetz makes the full kernel a finite torsion $W$-module. Every nonzero such module contains a nonzero element killed by two, so this kernel is zero. This proves the exact range.
\end{proof}


\section{Positive integral ranges}\label{sec:positive}

The examples concern small degrees for fixed embeddings. Sufficiently high degree gives an integral theorem. The following explicit regularity bound is a convenient form of the sufficiently ample argument of \citet[Lemma 2.12, pp.~242--243]{BMS}; it is not a claim that the high-degree theorem is new. Regularity below is Castelnuovo--Mumford regularity for the specified embedding, with coherent sheaves on the variety viewed as sheaves on its ambient projective space.

\begin{theorem}\label{thm:regularity}
Let $V\subset\PP^r_k$ be smooth projective of dimension $n\ge2$, and put
\[
 d_{\mathrm{int}}(V)=\max\bigl(1,\operatorname{reg}(\Omega_V^a):0\le a\le n\bigr).
\]
If $d\ge d_{\mathrm{int}}(V)$ and $Y\in|\OO_V(d)|$ is smooth, then crystalline restriction is an isomorphism below degree $n-1$ and injective with torsion-free cokernel in degree $n-1$. In that degree it identifies the full torsion subgroups. Coherent restriction in degree $n-1$ is injective already for
$d\ge\max(1,\operatorname{reg}\omega_V)$.
\end{theorem}
\begin{proof}
For every $t\ge d_{\mathrm{int}}(V)$, regularity gives
$H^b(V,\Omega_V^a(t))=0$ for $b>0$. By Serre duality and the wedge pairing,
\[
 H^b(V,\Omega_V^a(-t))=0\qquad(b<n).
\]
In particular this holds for $t=md$, $m\ge1$. Write $D^{\ge q}$ for complexes with zero cohomology in degrees below $q$. The divisor sequence then shows
\[
 R\Gamma(Y,\Omega_V^a|_Y(-md))\in D^{\ge n-1}\quad(m\ge1),
\]
and the cone of $R\Gamma(\Omega_V^a)\to R\Gamma(\Omega_V^a|_Y)$ is in $D^{\ge n-1}$.

For $0\le a\le n-1$, exterior powers of the conormal sequence give the exact resolution
\[
 0\to\OO_Y(-ad)\to\Omega_V^1|_Y(-(a-1)d)\to\cdots
 \to\Omega_V^{a-1}|_Y(-d)\to\Omega_V^a|_Y\to\Omega_Y^a\to0.
\]
Each of its $a$ preceding terms has cohomology in $D^{\ge n-1}$. Successive kernels therefore show
\[
 \operatorname{Cone}\bigl(R\Gamma(V,\Omega_V^a)
                   \to R\Gamma(Y,\Omega_Y^a)\bigr)
 \in D^{\ge n-a-1}.
\]
For $a=n$, the same assertion is immediate from $\Omega_Y^n=0$. Shifting each form-degree piece by $[-a]$ and using the finite de Rham filtration puts the de Rham restriction cone in $D^{\ge n-1}$.

Let $K$ be the crystalline restriction cone. It is a bounded complex with finite cohomology over $W$, and derived crystalline base change identifies $K\otimes_W^{\mathbf L}k$ with the de Rham cone. Its universal-coefficient sequences are
\[
 0\to H^j(K)/p\to H^j(K\otimes_W^{\mathbf L}k)
                 \to H^{j+1}(K)[p]\to0.
\]
Nakayama gives $H^j(K)=0$ for $j<n-1$. Taking $j=n-2$ also gives $H^{n-1}(K)[p]=0$, so that group is torsion-free. The long exact restriction sequence gives the lower-degree isomorphisms and the degree-$(n-1)$ injection. Its cokernel embeds into $H^{n-1}(K)$ and is therefore torsion-free. An injection of finite $W$-modules with torsion-free quotient identifies their torsion subgroups.

Finally, Serre duality gives
$H^{n-1}(V,\OO_V(-d))^\vee=H^1(V,\omega_V(d))$. The stated regularity bound kills this group, and the divisor sequence gives coherent injectivity.
\end{proof}

There are also familiar hypotheses giving the full range $d\ge1$. For a smooth global complete intersection $V\subset\PP^r$, every smooth hypersurface section has the integral conclusion of Theorem~\ref{thm:regularity}. The same holds if $n\le p$ and $V$ lifts smoothly to $W_2(k)$. Without that dimension bound, a $W_2$-lift gives isomorphism below degree $\min(n,p)-1$ and injection with torsion-free cokernel in that degree. These follow from the Hodge-equivalence statements of \citet[Definition 2.2.1 and Theorem 2.2.3, pp.~1854--1855]{CotnerZavyalov} by the de Rham filtration and universal-coefficient argument just proved. For complete intersections apply their comparison with projective space to $V$ and to $Y$ and use the triangle of the two cones. The lifting statement does not require a lift of the chosen divisor. None of these hypotheses is supplied by smoothness and very ampleness alone.


\section{Questions left by the examples}

The numerical size of the cover is not intended to be optimal. The following questions are distinct from the assertions proved above.

For the fixed odd-characteristic examples, determine exactly which degrees $2\le d<p^2-p$ have torsion-free second-crystalline restriction for every smooth member. The present calculation gives a sufficient vanishing range but does not determine these intermediate negative cohomology groups and relative de Rham maps.

Find a smaller explicit characteristic-two embedding with the same exact degree threshold. The simultaneous-root construction uses all 556 sections in order to recover the embedded base directly; replacing that basis by fewer sections requires a separate separation argument.

For the second threefold of Theorem~\ref{thm:injectivity-exact}, determine the full crystalline torsion groups and the exact degrees with torsion-free cokernel. The injectivity threshold alone does not determine either answer. Finding smaller embeddings with the same two distinct thresholds also remains open here.

\section*{Note on method and reuse}

The mathematics and exposition were produced by GPT-6 Astra Pro, the model name displayed by the interface, in ChatGPT Pro mode. Separate statement-only derivations by the same model cover the explicit surface, simultaneous-root cover, coherent degree threshold, differential bounds, product torsion growth, integral saturation criterion, every-characteristic extension, and iterated-Tango injectivity theorem. A targeted comparison identified an omitted restriction justification in the latter derivation; the same checker repaired it from the construction before seeing the printed proof. The original attempt and its correction are distinguished in the verification records. The written proofs and their published inputs also received targeted mathematical checks. This is not verification by an independent model family. No human has checked the paper. Exact arithmetic is reproduced by the programs \emph{Pole-semigroup certificate} and \emph{Iteration certificate} in the workbench at \url{\workbenchurl}. The permanent records are maintained through \url{\recordurl}. The text and original programs are dedicated to the public domain under CC0 1.0.

\begin{thebibliography}{99}
\bibitem[Bhatt et~al.(2018)]{BMS}
Bhatt, B., Morrow, M., \& Scholze, P. (2018). Integral $p$-adic Hodge theory. \emph{Publications mathématiques de l'IHES, 128}, 219--397. \url{https://doi.org/10.1007/s10240-019-00102-z}.
\bibitem[Cotner and Zavyalov(2024)]{CotnerZavyalov}
Cotner, S., \& Zavyalov, B. (2024). Lefschetz theorems in flat cohomology and applications. \emph{Compositio Mathematica, 160}(8), 1850--1870. \url{https://doi.org/10.1112/S0010437X24007231}.
\bibitem[Grothendieck(n.d.)]{Grothendieck}
Grothendieck, A. (n.d.). \emph{Torsion en cohomologie cristalline}. Archives Grothendieck, Université de Montpellier, cote 4, \S\S8.1--8.3, archive sheets 2--4 (PDF images 3--5). \url{https://grothendieck.umontpellier.fr/4.pdf}.
\bibitem[Hexagon(2026)]{Hexagon}
\emph{A short proof of Deligne's bound for Frobenius eigenvalues, after Alpöge}. (2026). Hexagon preprint 2611.00789v1. Contributor: F. Alderliefste. \url{https://hexagonmath.org/2611.00789v1}. The abstract and provenance state production from Alpöge's sketch using GPT-6 Astra; cited here for the stated scope, not as an input to a proof.
\bibitem[Illusie(1979)]{Illusie}
Illusie, L. (1979). Complexe de de Rham--Witt et cohomologie cristalline. \emph{Annales scientifiques de l'École normale supérieure, 12}(4), 501--661. \url{https://doi.org/10.24033/asens.1374}.
\bibitem[Mukai(2013)]{Mukai}
Mukai, S. (2013). Counterexamples to Kodaira's vanishing and Yau's inequality in positive characteristics. \emph{Kyoto Journal of Mathematics, 53}(2), 515--532. \url{https://doi.org/10.1215/21562261-2081279}. Page locators explicitly marked preprint refer to RIMS Preprint 1736 (2011), \url{https://www.kurims.kyoto-u.ac.jp/preprint/file/RIMS1736.pdf}.
\bibitem[Nakkajima and Yobuko(2022)]{NakkajimaYobuko}
Nakkajima, Y., \& Yobuko, F. (2022). \emph{Degenerations of log Hodge de Rham spectral sequences, log Kodaira vanishing theorem in characteristic $p>0$ and log weak Lefschetz conjecture for log crystalline cohomologies}. arXiv:1902.09110, version 3. Appendix by Y. Nakkajima, Theorem 9.1 and Corollary 9.3, pp.~63--64. \url{https://arxiv.org/abs/1902.09110}.
\bibitem[Raynaud(1978)]{Raynaud}
Raynaud, M. (1978). Contre-exemple au ``vanishing theorem'' en caractéristique $p>0$. In \emph{C. P. Ramanujam---A tribute} (pp.~273--278). Tata Institute of Fundamental Research / Springer.
\bibitem[Artin et~al.(1973)]{SGA4}
Artin, M., Grothendieck, A., \& Verdier, J.-L. (Eds.). (1973). \emph{Théorie des topos et cohomologie étale des schémas, tome 3} (SGA 4, Lecture Notes in Mathematics 305). Springer. \url{https://doi.org/10.1007/BFb0070714}.
\bibitem[Zheng(2017)]{Zheng}
Zheng, X. (2017). Counterexamples of Kodaira vanishing for smooth surfaces of general type in positive characteristic. \emph{Journal of Pure and Applied Algebra, 221}(10), 2431--2444. \url{https://doi.org/10.1016/j.jpaa.2016.12.030}. Statement and page locators in the text refer to arXiv:1509.04993v1, \url{https://arxiv.org/abs/1509.04993}.
\end{thebibliography}
\end{document}
